\section*{Capítulo \ref{ch:b2:current-potential-curves} --- Curvas intensidad--potencial}

\begin{solution}{exo:b2:current-potential-curves:1}
El cinc se oxida: \omterm{def:b2:current-potential-curves:ie-curve}{corriente anódica}, positiva. Los iones cobre se reducen en el
electrodo de cobre: \omterm{def:b2:current-potential-curves:ie-curve}{corriente catódica}, negativa. Las dos corrientes son iguales en
valor absoluto, pues la misma carga recorre el circuito.
\end{solution}

\begin{solution}{exo:b2:current-potential-curves:2}
$|i_{\lim}| = 96\,485 \times 0.50 \times 10^{-4} \times 1.6 \times 10^{-9} \times
2.0/(25 \times 10^{-6}) = \qty{6.2e-4}{A}$, es decir, \qty{0.62}{mA}.
\end{solution}

\begin{solution}{exo:b2:current-potential-curves:3}
Un \omterm{def:b2:current-potential-curves:fast-slow}{sistema rápido} corta el eje de potenciales bruscamente en su potencial de Nernst; un \omterm{def:b2:current-potential-curves:fast-slow}{sistema
lento} presenta a su alrededor un tramo plano de corriente nula, y sus ramas solo empiezan
tras una \omterm{def:b2:current-potential-curves:overpotential}{sobretensión}. Sobre platino: \ce{Fe^3+}/\ce{Fe^2+} (rápido);
\ce{O2}/\ce{H2O} (lento).
\end{solution}

\begin{solution}{exo:b2:current-potential-curves:4}
En el potencial de semionda, igual a $E^\circ = \qty{0.77}{V}$ cuando las dos formas
difunden igual.
\end{solution}

\begin{solution}{exo:b2:current-potential-curves:5}
Cerca de \qty{0.77}{V}, una onda catódica de \ce{Fe^3+} (rápida), que alcanza una meseta
proporcional a su concentración. Cerca de \qty{0.34}{V}, una segunda onda, al depositarse
cobre; su meseta se suma a la primera y es aproximadamente el doble de alta para una
concentración igual y coeficientes de difusión parecidos, ya que $n = 2$. Después, cerca de
\qty{0}{V} a pH 0 sobre platino (más abajo una vez que el electrodo se recubre de cobre),
el muro abrupto de la reducción de \ce{H+}.
\end{solution}

\begin{solution}{exo:b2:current-potential-curves:6}
A \qty{0.50}{V}: solo se reduce \ce{Fe^3+} a \ce{Fe^2+}, con su \omterm{def:b2:current-potential-curves:limiting-current}{corriente límite}.
A \qty{0.10}{V}: se reduce \ce{Fe^3+} y a la vez se deposita cobre,
cada uno con su \omterm{def:b2:current-potential-curves:limiting-current}{corriente límite}; el total es la suma de las dos mesetas.
\end{solution}

\begin{solution}{exo:b2:current-potential-curves:7}
pH 0: \qty{0}{V} y \qty{1.23}{V}; pH 14: \qty{-0.83}{V} y \qty{0.40}{V}. Sobre
platino la ventana conserva su anchura y baja \qty{59}{mV} por unidad de pH,
siguiendo cada muro a su par, y el anódico sigue desplazado hacia arriba por la
\omterm{def:b2:current-potential-curves:overpotential}{sobretensión} del oxígeno.
\end{solution}

\begin{solution}{exo:b2:current-potential-curves:8}
A \qty{1.0}{V}, el \ce{Fe^2+} se oxida en su meseta (\omterm{def:b2:current-potential-curves:ie-curve}{corriente anódica}
proporcional a $[\ce{Fe^2+}]$) y el \ce{Ce^4+} que haya se reduce en su meseta
(\omterm{def:b2:current-potential-curves:ie-curve}{corriente catódica} proporcional a $[\ce{Ce^4+}]$). Antes de la equivalencia la
\omterm{def:b2:current-potential-curves:ie-curve}{corriente anódica} disminuye linealmente con el volumen añadido; después, la \omterm{def:b2:current-potential-curves:ie-curve}{corriente
catódica} crece linealmente. Las dos rectas se anulan en la equivalencia.
\end{solution}

\begin{solution}{exo:b2:current-potential-curves:9}
Las mesetas se duplican ($|i_{\lim}| \propto 1/\delta$). El potencial de semionda no
cambia, porque la capa es la misma para ambas formas; tampoco cambia el potencial de corriente
nula de un \omterm{def:b2:current-potential-curves:fast-slow}{sistema rápido}, que es el potencial de Nernst.
\end{solution}

\begin{solution}{exo:b2:current-potential-curves:10}
Como en el capítulo, con $i_{\lim,c} = 0$: $E = E_{1/2} + (RT/nF)\ln[i/(i_{\lim} -
i)]$. En logaritmos decimales, $E = E_{1/2} + (RT\ln 10/nF)\log[i/(i_{\lim} - i)]$,
una recta de pendiente $0.0592/n$ voltios a \qty{298}{K}, que corta el
eje de $E$ en $E_{1/2}$.
\end{solution}

\begin{solution}{exo:b2:current-potential-curves:11}
Solo, el cinc se estabiliza en el potencial en el que su \omterm{def:b2:current-potential-curves:ie-curve}{corriente anódica} (oxidación
rápida del cinc) es igual a la \omterm{def:b2:current-potential-curves:ie-curve}{corriente catódica} de reducción de \ce{H+} sobre cinc,
que es pequeña porque esa reducción es lenta sobre cinc: la disolución es
lenta. En contacto con el platino, los electrones liberados por el cinc pueden reducir \ce{H+}
sobre el platino, donde la reducción es rápida: la \omterm{def:b2:current-potential-curves:ie-curve}{corriente catódica} a un potencial
dado es mucho mayor, el equilibrio de corrientes se alcanza a una corriente más alta, y el
cinc se disuelve más deprisa mientras se forma hidrógeno sobre el platino.
\end{solution}

\begin{solution}{exo:b2:current-potential-curves:12}
Sobre platino la oxidación del agua empieza cerca de \qty{1.6}{V} a pH 0: antes de
\qty{2.0}{V} la corriente se gastaría en producir dioxígeno. Un electrodo sobre el que
la oxidación del agua necesita una \omterm{def:b2:current-potential-curves:overpotential}{sobretensión} mayor desplaza el muro más allá de
\qty{2.0}{V}, de modo que la especie puede oxidarse dentro de la ventana.
\end{solution}

\begin{solution}{pb:b2:current-potential-curves:1}
\textbf{1.} \ce{O2 + 2H2O + 4e- -> 4OH-}.
\textbf{2.} \ce{Ag + Cl- -> AgCl + e-} (cuatro veces); global, \ce{O2 + 2H2O + 4Ag +
4Cl- -> 4AgCl + 4OH-}.
\textbf{3.} $E = 1.229 - 0.0592 \times 7 + (0.0592/4)\log 0.21 = \qty{0.80}{V}$.
\textbf{4.} $E^\circ(\ce{AgCl}/\ce{Ag}) = \qty{0.22}{V}$ con cloruro de actividad 1:
$-0.60 + 0.22 = \qty{-0.38}{V}$ en la escala del hidrógeno.
\textbf{5.} La reducción del \ce{O2} es un \omterm{def:b2:current-potential-curves:fast-slow}{sistema lento}: por debajo de su potencial de Nernst
hace falta una \omterm{def:b2:current-potential-curves:overpotential}{sobretensión} de varias décimas de voltio antes de que la corriente
crezca, y más aún para alcanzar la meseta.
\textbf{6.} En la meseta la corriente solo depende del aporte de \ce{O2},
no de las pequeñas derivas del potencial: mide la concentración.
\textbf{7.} La reducción del agua a dihidrógeno (el muro catódico), que
añadiría una corriente sin relación con el oxígeno.
\textbf{8.} Lineal a través de la membrana, de $c$ en el lado de la muestra a cero en
el cátodo.
\textbf{9.} $j = D_m c/L$ por unidad de área.
\textbf{10.} $i = -4FAD_m c/L$ (catódica), $|i| = 4FAD_mc/L$.
\textbf{11.} La membrana es la etapa más lenta: la difusión a su través es mucho
más lenta que en el agua exterior, de modo que la membrana desempeña el papel de la
\omterm{def:b2:current-potential-curves:limiting-current}{capa de difusión}.
\textbf{12.} El gradiente de concentración está dentro de la membrana, cuyo espesor
no depende de la agitación (siempre que el agua exterior se renueve);
un depósito sobre la membrana añade espesor y reduce la corriente.
\textbf{13.} $0.2095 \times (101.325 - 3.17) = \qty{20.56}{kPa}$.
\textbf{14.} $c = 1.3 \times 10^{-5} \times 20\,560 = \qty{0.27}{mol/m^3}$ ($0.267$).
\textbf{15.} $0.267 \times 32.00 = \qty{8.6}{mg/L}$ (\unit{g/m^3}).
\textbf{16.} $|i| = 4 \times 96\,485 \times 1.0 \times 10^{-5} \times 1.0 \times 10^{-10}
\times 0.267/(25 \times 10^{-6}) = \qty{4.1}{\micro A}$.
\textbf{17.} $4.05/8.55 = \qty{0.474}{\micro A}$ por \unit{mg/L}.
\textbf{18.} El espesor de la membrana y su coeficiente de difusión no se conocen
con precisión y cambian con la edad y la temperatura; una medida en una disolución
conocida los incluye a todos.
\textbf{19.} Cero (en la práctica, una pequeña corriente residual, que se resta).
\textbf{20.} $2.80/0.474 = \qty{5.9}{mg/L}$.
\textbf{21.} $2.80/4.05 = 69~\%$ de la saturación.
\textbf{22.} Oxígeno consumido por la descomposición de materia orgánica (aguas residuales, plantas muertas)
más deprisa de lo que lo aportan el aire y la fotosíntesis.
\textbf{23.} $2.80 \times 10^{-6}/(4 \times 96\,485) \times 3600 = \qty{2.6e-8}{mol}$
por hora: despreciable frente al oxígeno de un solo litro de agua
(\qty{1.8e-4}{mol}).
\textbf{24.} Tanto la solubilidad del \ce{O2} (\omterm{prop:b2:chemical-potential:henry}{constante de Henry}) como la difusión
a través de la membrana cambian con la temperatura: la sensibilidad hallada a
\qty{25}{\degreeCelsius} deja de ser válida.
\textbf{25.} La muestra contiene $\boldsymbol{\approx \qty{5.9}{mg/L}}$ de oxígeno
disuelto.
\end{solution}
